% JFMguide.tex  V1.1, 16 November 1993
% Copyright 1993 Cambridge University Press

% Cambridge University Press will substitute cm fonts for mt fonts
% (Monotype Times), so where you see \ifx\CUP@mtlplain@loaded\undefined
% this is an internal font switch. You may define new cm sizes etc.
% but if you do, please indicate with a `replace fonts!' message.
% Please keep the AMSTeX and font information together near the top
% of the file.

\documentstyle{JFM}

%\def\@mbi{mtmib10 at 10.5pt} % Does the style file need this def?

%%%%% USER-DEFINED MACROS HERE %%%%%
% Keep your own macro definitions separate from the examples below.
%
% Some useful examples of macro definitions follow.
%%%%%


%%%%% For units of measure %%%%%
%
\def\dynpercm{\nobreak\mbox{$\;$dynes\,cm$^{-1}$}}
\def\cmpermin{\nobreak\mbox{$\;$cm\,min$^{-1}$}}
%
%%%%%


%%%%% For multiletter symbols %%%%%
%
\def\Real{\mbox{Re}}      % cf plain TeX's \Re and Reynolds number
\def\Imag{\mbox{Im}}      % cf plain TeX's \Im
\def\Rey{\mbox{\it Re}}   % Reynolds number
\def\Pran{\mbox{\it Pr}}  % Prandtl number, cf plain TeX's \Pr product
\def\Pen{\mbox{\it Pe}}   % Peclet number
\def\Ai{\mbox{\rm Ai}}    % Airy function
\def\Bi{\mbox{\rm Bi}}    % Airy function
%
%%%%%


%%%%% For authors without AMSTeX %%%%%
% 
% Uncomment the definitions as far as: %%% STOP HERE without AMSTeX %%% 
%
\def\upi{\pi} % to be replaced with upright Greek character
\def\umu{\mu} % to be replaced with upright Greek character
\def\BbbE{\mbox{\sf E}} % to be replaced with blackboard bold 
\def\blackbox{\rule{4pt}{4pt}} % to be replaced with \blacksquare
%
%%%%% STOP HERE without AMSTeX %%%%%

\def\hexnumber#1{\ifcase#1 0\or1\or2\or3\or4\or5\or6\or7\or8\or9\or
 A\or B\or C\or D\or E\or F\fi }

%%%%% For authors with AMSTeX %%%%%
%
% Uncomment the definitions as far as %%% STOP HERE with AMSTeX %%% 
% 
% Example font defns for AMS Euler fonts for upright Greek characters
% 
% \font\eurmten=eurm10
% \font\eurmseven=eurm10 at 7pt
% \font\eurmfive=eurm10 at 5pt
% \newfam\eurmfam
% \textfont\eurmfam=\eurmten
% \scriptfont\eurmfam=\eurmseven
% \scriptscriptfont\eurmfam=\eurmfive
% \edef\eurm@{\hexnumber\eurmfam}
% 
% \mathchardef\upi="0\eurm@19   % for upright Greek character
% \mathchardef\umu="0\eurm@16   % for upright Greek character
%
% \font\msxten=msam10
% \font\msxseven=msam10 at 7pt
% \font\msxfive=msam10 at 5pt
% \newfam\msxfam
% \textfont\msxfam=\msxten
% \scriptfont\msxfam=\msxseven
% \scriptscriptfont\msxfam=\msxfive
% \edef\msx@{\hexnumber\msxfam}
%
% \mathchardef\leqslant="3\msx@36
% \mathchardef\geqslant="3\msx@3E
% \let\le=\leqslant
% \let\ge=\geqslant
% \let\leq=\leqslant
% \let\geq=\geqslant
%
% \def\BbbE{\Bbb E}
% \def\blackbox{\blacksquare} 
% \renewcommand{\simeq}{\approx}
%
%%%%% STOP HERE with AMSTeX %%%%%


%%%%% For sans serif characters %%%%%
%
\def\ssC{\mbox{\sf C}}         % for sans serif C
\def\sfsP{\mbox{\sfs P}}       % for sans serif sloping P
\def\slsQ{\mbox{\sls Q}}       % for sans serif bold-sloping Q
%
\makeatletter
\ifx\CUP@mtlplain@loaded\undefined
  \font\sfs = cmssi10    % sans-serif slanted
  \font\sfb = cmssi10    % sans-serif bold maths
  \font\sls = cmssi10    % sans-serif bold maths, slanted
  \font\bit = cmmib10    % bold math italic
\else
  % CUP times font substitutions for above (not to be used by authors)
  \font\sfs = mtssi10 at 10.5pt   % sans-serif italic
  \font\sfb = mtssbx10 at 10.5pt  % sans-serif bold maths
  \font\sls = mtssbi10 at 10.5pt  % sans-serif bold maths, slanted
  \font\bit = mtmib10 at 10.5pt   % bold math italic
\fi
\makeatother
%
%%%%%


%%%%% Hat position %%%%%
%
  \def\hatp{\skew5\hat{p}}        % p with hat
% \def\hatR{\skew5\hat{R}}        % R with hat 
% \def\hatRR{\skew5\hat{\hatR}}   % R with 2 hats  
% \def\doubletildesigma{\skew5\tilde{\skew5\tilde{\sSigma}}}     
%               italic Sigma with double tilde
%%%%%


%%%%% FOR MATH BOLD ITALIC %%%%%
%
% if you need math bold italic uncomment to END MATH BOLD ITALIC
%
 \makeatletter
 \ifx\CUP@mtlplain@loaded\undefined
   \font\tenbmi=cmmib10 at 10pt
   \font\sevenbmi=cmmib10 at 7pt
   \font\fivebmi=cmmib10 at 5pt

   \newfam\bmifam
   \textfont\bmifam=\tenbmi
   \scriptfont\bmifam=\sevenbmi
   \scriptscriptfont\bmifam=\fivebmi
   \def\bmi{\fam\bmifam\tenbmi}
 \fi
 \makeatother
%
%%%%% END MATH BOLD ITALIC %%%%%


% The use of \sbox, as shown here, is advised for frequently-occurring
% symbols, provided they always appear in the same font and point size.

\newsavebox{\thalfbox}
\sbox{\thalfbox}{$\textstyle\frac{1}{2}$}
\newcommand{\thalf}{\usebox{\thalfbox}}

\newsavebox{\shalfbox}
\sbox{\shalfbox}{$\scriptstyle\frac{1}{2}$}
\newcommand{\shalf}{\usebox{\shalfbox}}

\newsavebox{\squartbox}
\sbox{\squartbox}{$\frac{1}{4}$} %RM removed scriptstyle
\newcommand{\squart}{\usebox{\squartbox}}

\newsavebox{\etbox}
\sbox{\etbox}{\boldmath$\eta$}
%\sbox{\etbox}{$\eta$}  % use this if your \boldmath doesn't work.
\newcommand{\etb}{\usebox{\etbox}}

\newsavebox{\biSbox}
\sbox{\biSbox}{\bit S}
\newcommand{\biS}{\usebox{\biSbox}}

% array strut to make delimiters come out right size both ends
\newsavebox{\astrutbox}
\sbox{\astrutbox}{\rule[-5pt]{0pt}{20pt}}
\newcommand{\astrut}{\usebox{\astrutbox}}

% \newcommand{\seteqnum}[1]{\gdef\theequation{#1}}
% \newcommand{\reseteqnum}{\gdef\theequation{\thesection.\arabic{equation}}}

\ifnfsstwo
  \newcommand{\mitbf}[1] {\hbox{\mathversion{bold}${#1}$}}
  \newcommand{\rmn}[1] {{\mathrm #1}}
  \newcommand{\itl}[1] {{\mathit #1}}
  \newcommand{\bld}[1] {{\mathbf #1}}
\fi
\ifnfssone
  \newmathalphabet{\mathit}
    \addtoversion{normal}{\mathit}{cmr}{m}{it}
    \addtoversion{bold}{\mathit}{cmr}{bx}{it}
  \newmathalphabet{\mathcal}
    \addtoversion{normal}{\mathcal}{cmsy}{m}{n}
  \newcommand{\mitbf}[1] {\hbox{\mathversion{bold}${#1}$}}
  \newcommand{\rmn}[1] {{\mathrm #1}}
  \newcommand{\itl}[1] {{\mathit #1}}
  \newcommand{\bld}[1] {{\mathbf #1}}
\fi
\ifoldfss    
  \newcommand{\mitbf}[1] {\mbox{\boldmath $#1$}}
  \newcommand{\rmn}[1] {{\rm #1}}
  \newcommand{\itl}[1] {{\it #1}}
  \newcommand{\bld}[1] {{\bf #1}}
\fi

\newcommand{\GaPQ}  {G_a(P,Q)}
\newcommand{\Gat}   {\widetilde{G_a}}
\newcommand{\GsPQ}  {G_s(P,Q)}
\newcommand{\kgd}   {k\gamma d}
\newcommand{\ndq}   {\frac{\mbox{$\partial$}}{\mbox{$\partial$} n_q}}
\newcommand{\p}     {\mbox{$\partial$}}
\newcommand{\sh}    {^{\shalf}}
\newcommand{\smh}   {^{-\shalf}}
\newcommand{\sumjm} {\sum_{j=1}^{M}}
\newcommand{\tti}   {\rightarrow\infty}
\newcommand{\ttz}   {\rightarrow 0}
\mathchardef\varLambda="0103


%%%%% The minus sign is in different positions for the two fonts %%%%%
%
\makeatletter
\ifx\CUP@mtlplain@loaded\undefined
  \newcommand{\pvi}{\int_0^{\infty}\mskip -33mu -\quad} % 33mu for cm
\else
  \newcommand{\pvi}{\int_0^{\infty}\mskip -30mu -\quad} % 30mu for mt
\fi
\makeatother
%
%%%%%


%%%%% Bold Math Symbols %%%%%
%
\makeatletter
\ifx\CUP@mtlplain@loaded\undefined
  \font\tenbms=cmbsy10
  \font\sevenbms=cmbsy10 at 7pt
  \font\fivebms=cmbsy10 at 5pt
%
  \newfam\bmsfam
  \textfont\bmsfam=\tenbms
  \scriptfont\bmsfam=\sevenbms
  \scriptscriptfont\bmsfam=\fivebms
%
  \def\bms{\fam\bmsfam\tenbms}
  \edef\bsy@{\hexnumber\bmsfam}
  \mathchardef\bnabla="0\bsy@72
\fi
\makeatother
%
%%%%%

\def\eg{{e.g.\ }}
\def\etc{{etc.\ }}
\def\etal{\mbox{\it et al.}}

\newtheorem{lemma}{Lemma}
\newtheorem{corollary}{Corollary}

%-----title and author----------------------

\title[Journal of Fluid Mechanics]{\LaTeX\ Input Guide for Authors}

\author[\LaTeX\ Input Guide for Authors]{C\ls A\ls M\ls B\ls R\ls I\ls %
D\ls G\ls E\ns \TeX\ls -\ls T\ls O\ls -\ls T\ls Y\ls P\ls E}

\date{September 1993}
\pubyear{1993} % only needed when year of publication is not current year.
\volume{251}
\pagerange{\pageref{firstpage}--\pageref{lastpage}}

% -----------------------------

\begin{document}

\ifnfssone
\else
  \ifnfsstwo
  \else
    \ifoldfss
      \let\mathcal\cal
      \let\mathrm\rm
      \let\mathsf\sf
    \fi
  \fi
\fi

\label{firstpage}
\maketitle

\begin{abstract}
This guide is for authors who are preparing papers for the
{\em Journal of Fluid Mechanics}
using the \LaTeX\ document preparation system and the Cambridge
University Press JFM style file.
\end{abstract}

\tableofcontents

\section{Introduction}

The layout design for the {\em Journal of Fluid Mechanics} has been implemented as a \LaTeX\ style file.
The JFM style is based on the ARTICLE style as discussed in the \LaTeX\
manual.
Commands which differ from the standard \LaTeX\ interface, or which are
provided in addition to the standard interface, are explained in this
guide. This guide is not a substitute for the \LaTeX\ manual itself.

Note that the final printed version of papers will use a Monotype Times
font rather than the Computer Modern available to most authors. Also,
the measure in JFM style is different from \LaTeX\ article style. For
these reasons line lengths and page breaks will change and authors
should not insert hard breaks in their text. 

The sample pages (\verb"JFMsampl.tex") contain examples of most of the points
mentioned in this Guide, plus many other useful notes and definitions.

\subsection{Introduction to \LaTeX}

The \LaTeX\ document preparation system is a special version of the
\TeX\ typesetting program.
 \LaTeX\ adds to \TeX\ a collection of commands which allow the author
to concentrate on the logical structure of the document rather than its
visual layout.

\LaTeX\ provides a consistent and comprehensive document preparation
interface.
\LaTeX\ can automatically number equations, figures,
tables, and list entries, as well as sections and
subsections.
Using this numbering system, bibliographic citations, page references
and cross-references to any other numbered entity (\eg
section, equation, figure, list entry) are quite straightforward.

\subsection{The JFM document style}

The use of document styles allows a simple change of style (or style option) to
transform the appearance of your document.
The CUP JFM style file preserves the standard \LaTeX\ interface such that any
document which can be  produced using the standard \LaTeX\ ARTICLE style
can also be produced with the JFM style.
However, the measure (or width of text) is slightly different from that
for ARTICLE; therefore line breaks will change and it is possible that
equations may need re-setting.


\section{Using the JFM style file}

First, copy the file \verb"JFM.sty" into the correct subdirectory on your
system.
The JFM document style is implemented as a complete document style, {\em not\/}
a document style option.
In order to use the JFM style, replace \verb"article" by \verb"JFM" in the
\verb"\documentstyle" command at the beginning of your document:

\smallskip
\begin{verbatim}
  \documentstyle{article}
\end{verbatim}
\smallskip

\noindent
is replaced by

\smallskip
\begin{verbatim}
  \documentstyle{JFM}
\end{verbatim}
\smallskip

\noindent
None of the standard document style options should be used.
Author-defined macros should be inserted before \verb"\begin{document}",
or in a separate file and should be included with the submission,
see \S\ref{secUser}. Authors must not change any of the macro definitions
or parameters in \verb"JFM.sty".

\subsection{Document style options}

The standard document style options should {\em not\/} be used with
the JFM style:
\begin{itemize}
  \item {\tt 10pt}, {\tt 11pt}, {\tt 12pt} -- unavailable.
  \item {\tt draft}, {\tt twoside} (no associated style file) -- {\tt twoside}
        is the default.
  \item {\tt fleqn}, {\tt leqno}, {\tt titlepage}, {\tt twocolumn} --
        unavailable.
  \item {\tt proc}, {\tt ifthen}, {\tt bezier} -- can be used if necessary.
\end{itemize}

\section{Additional facilities}

In addition to all the standard \LaTeX\ design elements, the JFM style file
includes the following features:
\begin{itemize}
  \item Extended commands for specifying a short version
        of the title and author(s) for the running
        headlines.
  \item Abstract environment
  \item Control of enumerated lists.
  \item A \verb"proof" environment.
\end{itemize}
Once you have used these additional facilities in your document,
it can be processed only with \verb"JFM.sty".

\subsection{Titles, authors' names and running headlines}

In the JFM style, the title of the article and the author's name (or authors'
names) are used both at the beginning of the article for the main title and
throughout the article as running headlines at the top of every page.
The title is used on odd-numbered pages (rectos) and the author's name
(with initials only for first names) appears on even-numbered pages
(versos). The \verb"\pagestyle" and \verb"\thispagestyle" commands should
{\em not\/} be used.  Similarly, the commands \verb"\markright" and
\verb"\markboth" should not be necessary.

Although the main heading can run to several lines of text,
the running headline must be a single line.
Moreover, the main heading can also incorporate new line commands
(\eg \verb"\\") but these are not acceptable in a running headline.
To enable you to specify an alternative short title and author's name, the
standard \verb"\title" and \verb"\author" commands have been extended to take
an optional argument to be used as the running headline:

\smallskip
\begin{verbatim}
  \title[A short title]{The full title which can be as long
                         as necessary}
  \author[Initials and last names of all authors]{The full names of
     all the authors, using \and before the last name in the list}
  \affiliation{Author's affiliation}
\end{verbatim}
\smallskip

Unlike most other document styles, \verb"JFM.sty" does not follow the
convention of using the second line of the \verb"\author" command to
write the author's affiliation.  Rather, this must be entered using
the \verb"\affiliation" command as shown above.  The \verb"\and"
command does not produce separate author/affiliation pairs, but
only generates the word `and' in small caps as required
by the JFM style.  Use the predefined macro \verb"\ls" to get the
letterspacing between the letters of the author's name.  If there is
more than one affiliation, use \verb"\\[\affilskip]" between lines of the
\verb"\affiliation", inserting footnote numbers manually.

An example will make all this clearer.  To produce

\medskip
\begin{center}
{\Large\bf Interesting title} \par\medskip
{\bf By A\ls L\ls A\ls N\ns N.\ns O\ls T\ls H\ls E\ls R$^1$\ns
  \and J.\ns Q.\ns P\ls U\ls B\ls L\ls I\ls C$^2$}\par
\medskip
{\small $^1$University of Moscow, Moscow, Russia\\[3pt]
$^2$Ngonga University, Nairobi, Kenya}\par
\medskip
{\small (Received 12 June 1992)}
\end{center}
\medskip

\noindent you would type

\smallskip
\begin{verbatim}
  \title[Short title]{Interesting title}
  \author[A. N. Other and J. Q. Public]
         {A\ls L\ls A\ls N\ns N.\ns O\ls T\ls H\ls E\ls R$^1$\ns
          \and J.\ns Q.\ns P\ls U\ls B\ls L\ls I\ls C$^2$}
  \affiliation{$^1$University of Moscow, Moscow, Russia\\[\affilskip]
               $^2$Ngonga University, Nairobi, Kenya}
  \date{12 June 1992}
\end{verbatim}
\smallskip

The \LaTeX\ \verb"\thanks" command creates a footnote appearing at the bottom
of the title page.


\subsection{Abstract}

The JFM style provides for an abstract, produced by

\smallskip
\begin{verbatim}
  \begin{abstract}
  ...
  \end{abstract}
\end{verbatim}
\smallskip

This should appear just {\em before} the first \verb"\section" command.

\subsection{Lists}

The JFM style provides the three standard list environments plus an additional
unnumbered list:
\begin{itemize}
  \item Numbered lists, created using the \verb"enumerate" environment.
  \item Bulleted lists, created using the \verb"itemize" environment.
  \item Labelled lists, created using the \verb"description" environment.
\end{itemize}
The \verb"enumerate" environment numbers each list item with an italic
letter in parentheses;
alternative styles can be achieved by inserting a redefinition of the
number labelling command after the \verb"\begin{enumerate}". For example, a
list numbered with roman numerals inside parentheses can be produced by the
following commands:

\smallskip
\begin{verbatim}
  \begin{enumerate}
  \renewcommand{\theenumi}{\roman{enumi}}
  \item first item
         :
  \end{enumerate}
\end{verbatim}
\smallskip

\noindent
This produces the following list:
\begin{enumerate}
  \renewcommand{\theenumi}{\roman{enumi}}
  \item first item
  \item second item
  \item etc.
\end{enumerate}

In the last example, the item alignment is uneven because
the standard list indention is designed to be sufficient for arabic
numerals rather than the longer roman numerals. In order to enable
different labels to be used more easily, the \verb"enumerate" environment
in the JFM style can be given an optional argument which (like a standard
\verb"thebibliography" environment) specifies the {\em widest label}. For
example,
\begin{enumerate}[(iii)]
\renewcommand{\theenumi}{\roman{enumi}}
  \item first item
  \item second item
  \item etc.
\end{enumerate}
was produced by the following input:

\smallskip
\begin{verbatim}
  \begin{enumerate}[(iii)]
  \renewcommand{\theenumi}{\roman{enumi}}
  \item first item
          :
  \end{enumerate}
\end{verbatim}
\smallskip

Remember, once you have used the optional argument on the \verb"enumerate"
environment, do not process your document with a standard \LaTeX\ style file.




\subsection{Proof environment}

The standard \LaTeX\ constructs do not include a proof environment to
follow a theorem, lemma \etc (see also $\S$\ref{sectTheor}), and so one has
been added for the JFM style. For example,

\smallskip
\begin{verbatim}
  \begin{theorem}[2]
  Let the scalar function $T(x,y,t,{\bmi \omega}) $ be a conserved
  density for solutions of {\rm (9)}. Then the two-component function
  \begin{equation}
  {\mitbf{P}} = {\mathsf J}{\mathcal E} T
  \end{equation}
  represents the infinitesimal generator of a symmetry group
  for {\rm (9)}.
  \end{theorem}
  \begin{proof}
  The assumption about $T$ means that
  \[
  0 \sim \frac{{\mathrm \partial} T}{{\mathrm \partial} t} +
  {\mathcal E}T {\bmi \omega}_t
  = \frac{{\mathrm \partial} T}{{\mathrm \partial} t} + \{ T, H\},
  \]
  where ${\mathrm \partial} T / {\mathrm \partial} t$ refers to explicit
  dependence on $t$. The skew symmetry of  ${\mathsf J}$ hence implies
  \begin{equation}
  \frac{{\mathrm \partial} T}{{\mathrm \partial} t} \sim  \{ H, T\},
  \end{equation}
  whereupon the operation ${\mathsf J} \mathcal E$, which commutes with
  ${\mathrm \partial}_t$ in its present sense, gives
  \[
  \frac{{\mathrm \partial} {\mitbf{P}}}{{\mathrm \partial} t}
  = {\mathsf J}{\mathcal E}\{ H, T\}.
  \]
  This equation reproduces the characterisation of symmetries that was
  expressed by (19), thus showing ${\mathrm P}$ to represent a symmetry
  group.
  \end{proof}
\end{verbatim}
\smallskip

\noindent produces the following text:
   \begin{theorem}[2]
   Let the scalar function $T(x,y,t,{\bmi \omega}) $ be a conserved
   density for solutions of {\rm (9)}. Then the two-component function
   \begin{equation}
   {\mitbf{P}} = {\mathsf J}{\mathcal E} T
   \end{equation}
   represents the infinitesimal generator of a symmetry group for {\rm (9)}.
   \end{theorem}
   \begin{proof}
   The assumption about $T$ means that
   \[
   0 \sim \frac{{\mathrm \partial} T}{{\mathrm \partial} t} +
   {\mathcal E}T {\bmi \omega}_t
   = \frac{{\mathrm \partial} T}{{\mathrm \partial} t} + \{ T, H\},
   \]
   where ${\mathrm \partial} T / {\mathrm \partial} t$ refers to explicit
   dependence on $t$. The skew symmetry of  ${\mathsf J}$ hence implies
   \begin{equation}
   \frac{{\mathrm \partial} T}{{\mathrm \partial} t} \sim  \{ H, T\},
   \end{equation}
   whereupon the operation $\mathsf J \mathcal E$, which commutes with
   ${\mathrm \partial}_t$ in its present sense, gives
   \[
   \frac{{\mathrm \partial} {\mitbf{P}}}{{\mathrm \partial} t} = {\mathsf J}{\mathcal E}\{ H, T\}.
   \]
   This equation reproduces the characterisation of symmetries that was
   expressed by (19), thus showing ${\mathrm P}$ to represent a symmetry group.
   \end{proof}

The final \usebox{\proofbox} will not be included if the \verb"proof*" environment is
used.


\subsection{Remarks}

The JFM style requires Remarks, Notes, Problems, Comments, and Cases to be
typeset as for proofs but without the closing \usebox{\proofbox}.  The
\verb"JFM.sty" file contains a \verb"remarks" environment for this purpose.
It works like the standard \LaTeX\ \verb"\newtheorem" command except there
is no optional argument for numbering:

\smallskip
\begin{verbatim}
  \newremark{note}{Note}
  \begin{note}
     This is a note...
  \end{note}
\end{verbatim}
\smallskip

\noindent This will produce the following text:
\newremark{note}{Note}
\begin{note}
This is a note.  It isn't very interesting as notes go, but it's still
a note.
\end{note}

\section{Mathematics and units}

The JFM style file will insert the correct space above and below
displayed maths if standard \LaTeX\ commands are used; for example use
\verb"\[ ... \]" and {\em not\/} \verb"$$ ... $$". Do not leave blank
lines above and below displayed equations unless a new paragraph is
really intended.

\subsection{Numbering of equations}

\ifoldfss
The \verb"subequations" and \verb"subeqnarray" environments have been
incorporated into the JFM style file. Using these two environments, you can
number your equations (\ref{a1}), (\ref{a2}) \etc automatically. For
example, you can typeset
  \begin{subequations}
  \begin{equation}
    a_1 \equiv (2\Omega M^2/x)^{\frac{1}{4}} y^{\frac{1}{2}}\label{a1}
  \end{equation}
  and
  \begin{equation}
    a_2 \equiv (x/2\Omega)^{\frac{1}{2}}k_y/M.\label{a2}
  \end{equation}
  \end{subequations}
by using the \verb"subequations" environment as follows:
\smallskip
\begin{verbatim}
  \begin{subequations}
  \begin{equation}
    a_1 \equiv (2\Omega M^2/x)^{\frac{1}{4}} y^{\frac{1}{2}}\label{a1}
  \end{equation}
  and
  \begin{equation}
    a_2 \equiv (x/2\Omega)^{\frac{1}{2}}k_y/M.\label{a2}
  \end{equation}
  \end{subequations}
\end{verbatim}
\smallskip
You may also typeset an array such as:
  \begin{subeqnarray}
    \dot{X}    & = & \gamma X - \gamma\delta\eta ,\\
    \dot{\eta} & = & {\textstyle\frac{1}{2}} \delta + 2X\eta .
  \end{subeqnarray}
by using the \verb"subeqnarray" environment as follows:
\smallskip
  \begin{verbatim}
  \begin{subeqnarray}
    \dot{X}    & = & \gamma X - \gamma\delta\eta ,\\
    \dot{\eta} & = & {\textstyle\frac{1}{2}} \delta + 2X\eta .
  \end{subeqnarray}
  \end{verbatim} 
%
\else % if using the new font selection scheme
%
The \verb"subequations" and \verb"subeqnarray" environments have been
incorporated into the JFM style file, but they are only giving the correct
style (i.e.\ {\it a, b\/} \etc in the equation number) if you TeX your file
using the Old Font Selection Scheme. To get around this limitation, here is
one method by which you can typeset equations such as:
  \stepcounter{equation}
  $$ 
    a_1 \equiv (2\Omega M^2/x)^{\frac{1}{4}} y^{\frac{1}{2}}
    \eqno{(\theequation{\itl{a}})}
  $$
  and
  $$
    a_2 \equiv (x/2\Omega)^{\frac{1}{2}}k_y/M.\label{a2}
    \eqno{(\theequation{\itl{b}})}
  $$
You use the \verb"$$" notation (only in these exceptional circumstances):
\smallskip
\begin{verbatim}
  \stepcounter{equation}
  $$ 
    a_1 \equiv (2\Omega M^2/x)^{\frac{1}{4}} y^{\frac{1}{2}}
    \eqno{(\theequation{\itl{a}})}
  $$
  and
  $$
    a_2 \equiv (x/2\Omega)^{\frac{1}{2}}k_y/M.\label{a2}
    \eqno{(\theequation{\itl{b}})}
  $$
\end{verbatim}
\smallskip
You may also typeset an array such as:
  \newlength{\topleft}
  \newlength{\topright}
  \newlength{\bottomleft}
  \newlength{\bottomright}
  \newlength{\rightside} % difference in width on right side of equation
  \newlength{\leftside}  % difference in width on left side of equation
  \settowidth{\topleft}{$\dot{X}$}
  \settowidth{\topright}{$\gamma X - \gamma\delta\eta ,$}
  \settowidth{\bottomleft}{$\dot{\eta}$}
  \settowidth{\bottomright}{${\textstyle\frac{1}{2}} \delta + 2X\eta .$}
  \def\leftside{\topleft} \advance\leftside -\bottomleft
  \def\rightside{\topright} \advance\rightside -\bottomright
  \stepcounter{equation}
  $$
    \dot{X} = \gamma X - \gamma\delta\eta , \eqno{(\theequation{\itl{a}})}
  $$
  \vspace{-1.75\baselineskip} % this will vary!
  $$
    \hspace{\leftside}\dot{\eta} = {\textstyle\frac{1}{2}} \delta + 2X\eta .
    \hspace{\rightside}
    \eqno{(\theequation{\itl{b}})}
  $$
by using the following code; this simply makes the left- and right-hand sides
of both equations the same width, so that the equals signs align.
\smallskip
  \begin{verbatim}
  \newlength{\topleft}
  \newlength{\topright}
  \newlength{\bottomleft}
  \newlength{\bottomright}
  \newlength{\rightside} % difference in width on right side of equation
  \newlength{\leftside}  % difference in width on left side of equation
  \settowidth{\topleft}{$\dot{X}$}
  \settowidth{\topright}{$\gamma X - \gamma\delta\eta ,$}
  \settowidth{\bottomleft}{$\dot{\eta}$}
  \settowidth{\bottomright}{${\textstyle\frac{1}{2}} \delta + 2X\eta .$}
  \def\leftside{\topleft} \advance\leftside -\bottomleft
  \def\rightside{\topright} \advance\rightside -\bottomright
  \stepcounter{equation}
  $$
    \dot{X} = \gamma X - \gamma\delta\eta , \eqno{(\theequation{\itl{a}})}
  $$
  \vspace{-1.75\baselineskip} % this will vary!
  $$
    \hspace{\leftside}\dot{\eta} = {\textstyle\frac{1}{2}} \delta + 2X\eta .
    \hspace{\rightside}
    \eqno{(\theequation{\itl{b}})}
  $$
  \end{verbatim} 
\fi % end of NFSS version


\subsection{AMS fonts -- especially if you do not have them}

If you need AMS symbols but do not have the AMS fonts you can ensure that
they will be typeset in your article by taking the following steps. Set
up user-defined macros that can be redefined by the typesetter to use
AMS\TeX\ macros. For example, the blackboard bold symbols, sometimes
called shell or outline characters, are obtained with the AMS\TeX\ command
\verb"{\Bbb E}". Instead, use a macro definition such as:

\smallskip
\begin{verbatim}
   % replace font!
   \def\BbbE{\mbox{\sf E}}   % sans serif E
\end{verbatim}
\smallskip

\noindent This substitutes a sans serif character where you want blackboard bold.
You can typeset the input file and the typesetter is alerted to do the
substitution.

Plain \TeX\ provides only \verb"\geq" and \verb"\leq" which typeset the
Computer Modern symbols $\geq$ and $\leq$, respectively. These will be
redefined at typesetting to use the AMS\TeX\ macros \verb"\geqslant" and
\verb"\leqslant", to give the slanted symbols from the AMS fonts.

If you are using AMS fonts make sure they are {\em at least} version 2.1.
An earlier version of AMS fonts will cause problems in typesetting. You
should submit your copy of \verb"mssymb.tex" along with your input file.
If, in addition to the \verb"msxm" and \verb"msym" symbol fonts, you want
to use the AMS Euler fonts you should include the type of font definitions
present in the preamble of this guide's input file.

\subsection{Typeface}

Sometimes, non-italic symbols are required in maths, as described in the
Notation Guide (\verb"JFMnot.tex"). This section describes how these may be obtained using \LaTeX\ and \verb"JFM.sty".

\subsubsection{Roman symbols}\label{roman}
The mathematical operators and constants, such as sin, cos, log and exp,
are covered by plain \TeX\ macros which ensure that they are typeset in
roman text, even in math mode: \verb"\sin", \verb"\cos", \verb"\log",
\verb"\exp". Where single letters are concerned (\eg d, i, e) just use
the roman typeface in maths, i.e.\ \verb"$E=m{\mathrm c}^2$" which typesets as
$E=m{\mathrm c}^2$, giving the correct roman character but with maths spacing.
When the term involves more than one character (\eg \Real\ or \Imag )
text-character spacing is required:

\smallskip
\begin{verbatim}
  $\mbox{Re}\;x$
\end{verbatim}
\smallskip

\noindent which typesets as $\Real\;x$.

Where such expressions are used repeatedly, macro definitions can reduce
typing and editing. The following examples are included in the preamble
of the input files for this document, \verb"JFMguide.tex", and of the
sample article \verb"JFMsampl.tex". Authors are encouraged to use them
and others like them.

\smallskip
\begin{verbatim}  
  \def\Real{\mbox{Re}}   % do not confuse with TeX's \Re
  \def\Imag{\mbox{Im}}   % do not confuse with TeX's \Im
  \def\Ai{\mbox{Ai}}     % Airy function
  \def\Bi{\mbox{Bi}}     % Airy function
\end{verbatim}

\subsubsection{Multiletter italic symbols}
If multiletter symbols are used in maths mode, for example Reynolds, 
Prandtl numbers, \etc the standard maths mode spacing between them is 
too large and text-character spacing is required. As described 
in \S\ref{roman} (but here for italic letters) use for example

\smallskip
\begin{verbatim}
  \def\Rey{\mbox{\it Re}}
  \def\Pran{\mbox{\it Pr}}
\end{verbatim}

\subsubsection{Upright and sloping Greek symbols}
Like the italic/roman fonts described above, Greek letters for maths
variables are printed in the Journal in sloping type, whereas constants,
operators and units are upright. However, plain TeX macros produce sloping
lowercase Greek and upright capitals. The upright lowercase required for
example for `pi' and `mu' (micro)  are available in the AMS Euler fonts
(version 2.1) but you may not have
these. In this case, you will have to use the normal math italic symbols
and the typesetter will substitute the corresponding upright characters.
You will make this much easier if you can use the macros \verb"\upi" and
\verb"\umu" in your text to indicate the need for the upright characters,
together with the following definitions in the preamble, where they can
easily be redefined by the typesetter.

\smallskip
\begin{verbatim}
  \def\upi{\pi}
  \def\umu{\mu}
\end{verbatim}
\smallskip

Consult AMS\TeX\ for the sloping, uppercase Greek characters
(\verb"\varGamma" etc.). The few upright capital letters required 
(e.g.\ cap.\ delta for `difference') are given by the relevant plain 
TeX macro.

\subsubsection{Sans serif symbols}
The plain \TeX\ command \verb"\sf" changes the typeface to sans serif,
giving upright characters. Occasionally, bold-sloping sans serifs are
needed (see the Notation Guide). You should use the following types of
macros to obtain these fonts.

\smallskip
\begin{verbatim}
  \def\ssC{\mbox{\sf C}}     % for sans serif C
  \def\sfsP{\mbox{\sfs P}}   % for sans serif slanted P
  \def\sfbsX{\mbox{\sls X}}  % for sans serif bold slanted X
     .
     .
  % replace fonts!

  \font\sfs = cmssi10    % sans-serif slanted
  \font\sfb = cmssi10    % sans-serif bold maths 
  \font\sls = cmssi10    % sans-serif bold maths, slanted 
\end{verbatim}
\smallskip

\noindent Note that the bold-slanted font \verb"sls" uses the slanted sans
serif \verb"cmssi" -- because there is no bold-slanted maths sans serif
font in Computer Modern! If you can follow this procedure, the
typesetter will be able to substitute the proper bold-slanted symbol,
even though the proofs that you produce in Computer Modern will only use
the slanted font.

\subsubsection{Script characters}
Script characters should be typeset using plain \TeX 's \verb"\mathcal"
command. This produces the Computer Modern symbols such as {$\mathcal E\,$}
and {$\mathcal F\,$} in your hard copy but the the typesetter will substitute
the more florid script characters normally seen in the journal.

\subsection{Skewing of accents}
Accents such as hats, overbars and dots are normally centred over letters,
but when these are italic or sloping greek the accent may need to be moved
to the right so that it is centred over the top of the sloped letter. For
example, \verb"\def\hatp{\skew5\hat{p}}" will produce $\hatp$.

\subsection{Units of measure}
Numbers and their units of measure should be typeset with fixed spaces
that will not break over two lines. This is easily done with user-defined
macros. For example,\linebreak \verb"52\dynpercm" typesets as
52\dynpercm , providing the following macro definition has been included
in the preamble.

\smallskip
\begin{verbatim}
  \def\dynpercm{\nobreak\mbox{$\;$dynes\,cm$^{-1}$}}
\end{verbatim}

\section{User-defined macros}\label{secUser}
If you define your own macros you must ensure that their names do not
conflict with any existing macros in plain \TeX\ or \LaTeX\ (or AMS\TeX\ 
if you are using this). You should also place them in the preamble to
your input file, between the \verb"\documentstyle" and
\verb"\begin{document}" commands.

Apart from scanning the indexes of the relevant manuals, you can check
whether a macro name is already used in plain \TeX\ or \LaTeX\ by using
the \TeX\ command \verb"show". For instance, run \LaTeX\ interactively
and type \verb"\show\<macro_name>" at the \TeX\ prompt. (Alternatively,
insert the \verb"\show" command into the preamble of an input file and
\TeX\ it.)

\smallskip
\begin{verbatim}
  * \show\Re
\end{verbatim}
\smallskip

\noindent produces the following response from \TeX :

\smallskip
\begin{verbatim}
  > \Re=\mathchar"23C.
  <*> \show\Re
\end{verbatim}
\smallskip

\noindent By contrast \verb"\Real" is not part of plain \TeX\ or \LaTeX\ and
\verb"\show\Real" generates:

\smallskip
\begin{verbatim}
  > \Real=undefined
  <*> \show\Real
\end{verbatim}
\smallskip

\noindent confirming that this name can be assigned to a user-defined macro.

Such macros must be in a place where they can easily be found and
modified by the journal's editors or typesetter. They must be gathered
together in the preamble of your input file, or in a separate
\verb"macros.tex" file with the command \verb"\include{macros}" in the
preamble. Macro definitions must not be scattered about your document
where they are likely to be completely overlooked by the typesetter.

The same applies to font definitions that are based on Computer Modern
fonts. These must be changed by the typesetter to use the journal's
\TeX\ typefaces Times and Helvetica. In this case, you should draw
attention to these font definitions on the hard copy that you submit for
publication and by placing a comment in your input file just before the
relevant definitions, for example \verb"% replace font!"


\section{Some guidelines for using standard facilities}

The following notes may help you achieve the best effects with the
standard \LaTeX\ facilities that remain in the JFM style.

\subsection{Sections}

Only the first three \LaTeX\ section levels are defined in the JFM style
file:
\begin{itemize}
  \item[] Heading A -- \verb"\begin{section}"
  \item[] Heading B -- \verb"\begin{subsection}"
  \item[] Heading C -- \verb"\begin{subsubsection}"
\end{itemize}
There is no \verb"paragraph" or \verb"subparagraph" in the JFM style.

To obtain non-bold in a bold heading use the usual plain \TeX\ commands for
changing typeface; for example \verb"\section{Fluctuations in Ca{\sc ii}}".

\begin{table}
  \begin{center}
   \begin{minipage}{6cm}
    \begin{tabular}{cccc}
     {Figure} & {$hA$} & {$hB$}\footnote{A table must be 
      inside a \verb"minipage" environment if it includes 
      table footnotes.}
       & {$hC$}\\
     2 & $\exp\;(\upi {\mathrm i} x)$
       & $\exp\;(\upi {\mathrm i} y)$ & $0$\\
     3 & $-1$    & $\exp\;(\upi {\mathrm i} x)$ & $1$\\
     4 & $-4+3{\mathrm i}$ & $-4+3{\mathrm i}$ & 1.6\\
     5 & $-2$    & $-2$    & $1.2 {\mathrm i}$
    \end{tabular}
   \end{minipage}
  \end{center}
\caption{An example table} \label{sample-table}
\end{table}
%
\subsection{Tables}

The {\tt table} environment is implemented as described in
the \LaTeX\ manual to provide consecutively numbered floating 
inserts for tables.

JFM style will cope with most table positioning problems
and you should not normally use the optional positional qualifiers \verb"t",
\verb"b", \verb"h" on the
\verb"table" environment, as this would override these decisions.
Table captions should appear at the bottom of the table; therefore you
should place the \verb"\caption" command after the body of the table.

The JFM \verb"\table" command will insert rules above and below
the table as required by the JFM style, so you should not attempt 
to insert these yourself. 
The only time you need to intervene with the rules is when two or more tables
fall one above another. When this happens, you will get two rules together
and too much space between the tables; the solution is to add \verb"\norule"
before the end of upper table and \verb"\followon" between tables as shown here:

\smallskip
\begin{verbatim}
    ...
    \end{center}
    \caption{Dimensionless parameters}\label{dimensionless_p}
    \norule  
  \end{table}
    \followon
  \begin{table} 
    \begin{center}
    ... 
\end{verbatim}
\smallskip

The JFM style dictates that vertical rules should never be used within the 
body of the table, and horizontal rules should be used only to span 
columns with the same headings. Extra space can be inserted to distinguish
groups of rows or columns.

As an example, table~\ref{sample-table} is produced using the following
commands:

\smallskip
\begin{verbatim}
  \begin{table}
    \begin{center}
     \begin{minipage}{6cm}
      \begin{tabular}{cccc}
       {Figure} & {$hA$} & {$hB$}\footnote{A table must be 
        inside a \verb"minipage" environment if it includes 
        table footnotes.}
         & {$hC$}\\
       2 & $\exp\;(\upi {\mathrm i} x)$
         & $\exp\;(\upi {\mathrm i} y)$ & $0$\\
       3 & $-1$    & $\exp\;(\upi {\mathrm i} x)$ & $1$\\
       4 & $-4+3{\mathrm i}$ & $-4+3{\mathrm i}$ & 1.6\\
       5 & $-2$    & $-2$    & $1.2 {\mathrm i}$
      \end{tabular}
     \end{minipage}
    \end{center}
  \caption{An example table} \label{sample-table}
  \end{table}
\end{verbatim}
\smallskip

The \verb"tabular" environment
has been modified for the JFM style in the following ways:
\begin{enumerate}
  \item Additional vertical space is inserted above and below a horizontal   
        rule produced by \verb"\hline"
  \item Tables are centred, and span the full width of the page; that is,
        they are similar to the tables that would be produced by
        \verb"\begin{tabular*}{\textwidth}".
\end{enumerate}
Commands to
redefine quantities such as \verb"\arraystretch" should be omitted. If
the old tabular facilities are needed, there is a new environment,
\verb"oldtabular", which has none of the reformatting; it should be used
in exactly the same way.

\subsection{Illustrations (or figures)}

At present, figures cannot be handled using PostScript or similar files 
and so artwork must be supplied separately as hard copy to be lettered
and sized by the Printer. An approximate amount of space should be left, using the \verb"\vspace" command.

The JFM style will cope with most figure positioning problems and you should not
normally use the optional positional qualifiers \verb"t", \verb"b", \verb"h"
on the \verb"figure" environment, as this would override these decisions.
Figure captions should be below the figure itself, therefore the \verb"\caption"
command should appear after the space left for the illustration within the
\verb"figure" environment.
For example, figure~\ref{sample-figure} is produced using the following
commands:
\begin{figure}
  \vspace{3cm}
  \caption{An example figure with space for artwork}
  \label{sample-figure}
\end{figure}

\smallskip
\begin{verbatim}
  \begin{figure}
    \vspace{3cm}
    \caption{An example figure with space for artwork}
    \label{sample-figure}
  \end{figure}
\end{verbatim}
\smallskip

\subsection{Theorem environments}\label{sectTheor}

The \verb"\newtheorem" command works as described in the \LaTeX\ manual,
but produces spacing and caption typefaces required to the JFM style.  The
preferred numbering scheme is for theorems to be numbered within sections,
as 1.1, 1.2, 1.3, etc., but other numbering schemes are permissible and
may be implemented as described in the \LaTeX\ manual.  In order to allow
authors maximum flexibility in numbering and naming, {\em no} theorem-like
environments are defined in \verb"JFM.sty".  Rather, you have to define
each one yourself.  Theorem-like environments include Theorem, Definition,
Lemma, Corollary, and Proposition.  For example,

\smallskip
\begin{verbatim}
  \newtheorem{lemma}{Lemma}[section]
  \renewcommand{\thelemma}{\Roman{lemma}}
\end{verbatim}
\smallskip


\subsection{Acknowledgments}

Acknowledgments should appear at the close of your paper, just before
the list of references and any appendices.  Use the \verb"acknowledgments"
environment, which simply inserts an appropriate amount of vertical space.

\subsection{Appendices}

You should use the standard \LaTeX\ \verb"\appendix" command to place any
Appendices, normally, just before the references. This numbers
appendices as A, B \etc, equations as (A1), (B1) \etc, and figures and
tables as A1, B1 \etc

\subsection{References}

As with standard \LaTeX, there are two ways of producing a list of
references; either by compiling a list (using a \verb"thebibliography"
environment), or by using BibTeX with a suitable bibliographic database.

\subsubsection{References in the text}

References in the text are given by author and date, in the form
\cite{Den85}. Each  entry has a key, which is assigned by the author 
and used to refer to that entry in the text.

\subsubsection{The list of references}

The following listing shows some references prepared in the style of the
journal; the code produces the references at the end of this guide.

\smallskip
\begin{verbatim}
\begin{thebibliography}{} 
  \bibitem[Abramowitz \& Stegun (1965)]{AS65}
    {\sc Abramowitz, M. \& Stegun, I. A.} 1965
    {\em Handbook of Mathematical Functions}. Dover.
  \bibitem[Dennis (1985)]{Den85}
    {\sc Dennis, S. C. R.} 1985
    Compact explicit finite difference approximations to the
    Navier--Stokes equation. In {\em Ninth Intl Conf.\ on
    Numerical Methods in Fluid Dynamics\/} (ed.\ Soubbaramayer
    \& J. P. Boujot). Lecture Notes in Physics, vol.\ 218,
    pp.~23--51. Springer.
  \bibitem[Jones (1976)]{Jon76}
    {\sc Jones, O. C.} 1976
    An improvement in the calculation of turbulent friction in
    rectangular ducts. {\em Trans.\ ASME\/} J:
    {\em J. Fluids Engng.\ \/}{\bf 98}, 173--181.
  \bibitem[Saffman (1990)]{Saf90}
    {\sc Saffman, P. G.} 1990
    A model vortex reconnection. {\em J. Fluid Mech.\/}
    {\bf 212}, 395--402.
  \bibitem[Saffman \& Schatzman (1982)]{SS82}
    {\sc Saffman, P. G. \& Schatzman, J. C.} 1982
    Stability of a vortex street of finite vortices.
    {\em J. Fluid Mech.\ \/}{\bf 117}, 171--185.
  \bibitem[Saffman \& Yuen (1980)]{SY80}
    {\sc Saffman, P. G. \& Yuen, H. C.} 1980
    A new type of three-dimensional deep-water wave of permanent
    form. {\em J. Fluid Mech.\ } {\bf 101}, 797--808.
  \bibitem[Shaqfeh \& Koch (1990)]{SK90}
    {\sc Shaqfeh, E. S. G. \& Koch, D. L.} 1990
    Orientational dispersion of fibres in extensional flow.
    {\em Phys.\ Fluids} A {\bf 2}, 1077--1081.
  \bibitem[Wijngaarden (1968)]{Wij68}
    {\sc Wijngaarden, L. van} 1968
    On the oscillations near and at resonance in open pipes.
    {\em J. Engng Maths\/} {\bf 2}, 225--240
  \bibitem[Williams (1964)]{Wil64}
    {\sc Williams, J. A.} 1964
    A nonlinear problem in surface water waves. PhD thesis,
    University of California, Berkeley.
\end{thebibliography}
\end{verbatim}
\smallskip

\noindent
Each entry takes the form

\smallskip
\begin{verbatim}
  \bibitem[Author(s) (Date)]{tag}
    Bibliography entry
\end{verbatim}
\smallskip

\noindent
where \verb"Author(s)"\ should be the author names as they are cited in
the text, \verb"Date" is the date to be cited in the text, and \verb"tag"
is the tag that is to be used as an argument for the \verb"\cite{}" command.
\verb"Bibliography entry" should be the material that is to appear in the
bibliography, suitably formatted.

\begin{thebibliography}{} 
  \bibitem[Abramowitz \& Stegun (1965)]{AS65}
    {\sc Abramowitz, M. \& Stegun, I. A.} 1965
    {\em Handbook of Mathematical Functions}. Dover.
  \bibitem[Dennis (1985)]{Den85}
    {\sc Dennis, S. C. R.} 1985
    Compact explicit finite difference approximations to the
    Navier--Stokes equation. In {\em Ninth Intl Conf.\ on
    Numerical Methods in Fluid Dynamics\/} (ed.\ Soubbaramayer
    \& J. P. Boujot). Lecture Notes in Physics, vol.\ 218,
    pp.~23--51. Springer.
  \bibitem[Jones (1976)]{Jon76}
    {\sc Jones, O. C.} 1976
    An improvement in the calculation of turbulent friction in
    rectangular ducts. {\em Trans.\ ASME\/} J:
    {\em J. Fluids Engng.\ \/}{\bf 98}, 173--181.
  \bibitem[Saffman (1990)]{Saf90}
    {\sc Saffman, P. G.} 1990
    A model vortex reconnection. {\em J. Fluid Mech.\/}
    {\bf 212}, 395--402.
  \bibitem[Saffman \& Schatzman (1982)]{SS82}
    {\sc Saffman, P. G. \& Schatzman, J. C.} 1982
    Stability of a vortex street of finite vortices.
    {\em J. Fluid Mech.\ \/}{\bf 117}, 171--185.
  \bibitem[Saffman \& Yuen (1980)]{SY80}
    {\sc Saffman, P. G. \& Yuen, H. C.} 1980
    A new type of three-dimensional deep-water wave of permanent
    form. {\em J. Fluid Mech.\ } {\bf 101}, 797--808.
  \bibitem[Shaqfeh \& Koch (1990)]{SK90}
    {\sc Shaqfeh, E. S. G. \& Koch, D. L.} 1990
    Orientational dispersion of fibres in extensional flow.
    {\em Phys.\ Fluids} A {\bf 2}, 1077--1081.
  \bibitem[Wijngaarden (1968)]{Wij68}
    {\sc Wijngaarden, L. van} 1968
    On the oscillations near and at resonance in open pipes.
    {\em J. Engng Maths\/} {\bf 2}, 225--240
  \bibitem[Williams (1964)]{Wil64}
    {\sc Williams, J. A.} 1964
    A nonlinear problem in surface water waves. PhD thesis,
    University of California, Berkeley.
\end{thebibliography}
\newpage

\appendix
\section{Special commands in JFM.sty}

The following is a summary of the new commands, optional
arguments and environments which have been added to the
standard LaTeX user-interface in creating the JFM style file.

\vspace{6pt}

\begin{tabular}{lp{8cm}}
{\em New commands}    & \\ 
\verb"\affiliation"   & use after \verb"\author" to typeset an affiliation.\\
\verb"\affilskip"     & use immediately after \verb"\\" in
                        \verb"\affiliation" to
                        give correct vertical spacing between separate
                        affiliations. \\
\verb"\author"        & do not use \verb"\\" in \verb"\author" to
                        start an affiliation. \\
\verb"\followon"      & to remove space between adjacent tables.\\
\verb"\ls"            & to letter space the author's name. \\
\verb"\newremark"     & used as \verb"\newtheorem" to define environments
                        for Remarks, Definitions, Notes, etc. \\
\verb"\norule"        & to remove the line below a table.\\
\verb"\ns"            & to insert space between an author's names.  
                                                          \\[6.5pt]
{\em New optional arguments} & \\
\verb"[<short title>]" & in the \verb"\title" command: to define a right
                         running headline that is different from the article
                         title. \\
\verb"[<short author>]" & in the \verb"\author" command: to define a left
                          running headline that is different from the
                          authors' names as typeset at the article opening.\\
\verb"[<widest label>]" & in \verb"\begin{enumerate}": to ensure the correct
                          alignment of numbered lists.\\[6.5pt]
{\em New environments}  & \\
\verb"proof"            & to typeset mathematical proofs. \\
\verb"proof*"           & to typeset mathematical proofs without the
                          terminating proofbox. \\
\verb"remark"           & this environment works like the \verb"theorem"
                          environment;
                          it typesets an italic heading and roman text to
                          contrast with \verb"theorem"'s small caps heading
                          and italic text.\\
\verb"subeqnarray"      & enables equation numbers in an array to be 
                          numbered as (6.1a), (6.2b), etc.\\
\verb"subequation"      & enables consecutive equations to be numbered
                          (6.1a), (6.1b), etc.\\
\verb"tabular"          & has been modified to insert additional space above  
                          and below an \verb"\hrule" and the table caption
                          and body is centred with rules full out across the
                          text measure.
\end{tabular}

\newpage

\section{Notes for editors}
\setcounter{subsubsection}{0}

This appendix contains additional information which may be useful to
those who are involved with the final production stages of an article.
Authors, who are generally not typesetting the final pages in the
journal's typeface (Times), do not need this information.

\subsubsection{Catchline commands}
To be placed in the preamble:
\begin{itemize}
  \item \verb"\pubyear{}"
  \item \verb"\volume{}"
  \item \verb"\pagerange{}"
\end{itemize}

\subsubsection{Footnotes}
If a footnote falls at the bottom of a page, it is possible for the
footnote to appear on the following page (a feature of \TeX ). Check
for this.

\subsubsection{Rules between tables}
If two or more tables fall one above another, add the commands 
\verb"\norule" and \verb"\followon" as described in this guide.

\subsubsection{Font substitution}
Check for use of AMS fonts, bold slanted sans serifs, and bold math italic
and alter preamble definitions to use the appropriate AMS/CUP/Monotype
fonts for phototypesetter output.

\subsubsection{Font sizes.}
The JFM style file defines all the standard \LaTeX\ font sizes. For example:
\begin{itemize}
  \item \verb"tiny" -- {\tiny This is tiny text.}
  \item \verb"scriptsize" -- {\scriptsize This is scriptsize text.}
  \item \verb"\indexsize" -- {\indexsize This is indexsize text.}
  \item \verb"\footnotesize" -- {\footnotesize This is footnotesize text.}
  \item \verb"\small" -- {\small This is small text.}
  \item \verb"\qsmall"  -- {\qsmall This is qsmall text (quotations).}
  \item \verb"\normalsize" -- This is normalsize text (default).
  \item \verb"\large" -- {\large This is large text.}
  \item \verb"\Large" -- {\Large This is Large text.}
  \item \verb"\LARGE" -- {\LARGE This is LARGE text.}
\end{itemize}
%
All these sizes are summarized in Table~\ref{tab:fontsizes}.
%
\begin{table}
 \begin{tabular}{lr@{/}lp{4cm}}
 \multicolumn{1}{l}{\it Size} &
  \multicolumn{2}{c}{\it size/baseline} & \it Usage\\
 \hline
 \verb"\tiny"         &     5pt &  6pt    & -- \\
 \verb"\scriptsize"   &     7pt &  8pt    & -- \\
 \verb"\indexsize"    &     8pt &  9pt    & catchline. \\
 \verb"\footnotesize" &     9pt & 10pt    & index, footnotes. \\
 \verb"\small"        &     9pt & 10pt    & `AND' in authors' names,
                                            received date, affiliation,
                                            figure and table captions. \\
 \verb"\qsmall"       &  9.75pt & 10.75pt & quote, quotations. \\
 \verb"\normalsize"   &    10pt & 12pt    & main text size, abstract,
                                            B~and~C heads. \\
 \verb"\large"        &    11pt & 13pt    & A head, part no. \\
 \verb"\Large"        &    14pt & 18pt    & part title. \\
 \verb"\LARGE"        &    17pt & 19pt    & article title. \\
 \verb"\huge"         &    20pt & 25pt    & -- \\
 \verb"\Huge"         &    25pt & 30pt    & -- \\
 \end{tabular}
\caption{Type sizes for \LaTeX\ size-changing commands.}
\label{tab:fontsizes}
\end{table}


\label{lastpage}

\end{document}
